\documentclass[12pt]{article}
\usepackage[utf8]{inputenc}
\usepackage[T1]{fontenc}
\usepackage{hyperref}
\title{Tai Thanh Language Resources for K-12 Curriculum Development}
\author{Compiled by Claude Sonnet 5 (claude-sonnet-5) for Andr\'as Kornai}
\date{}
\begin{document}
\maketitle

\section{Introduction}
Tai Thanh (ISO 639-3: tmm) is a Southwestern Tai language of the Kra-Dai family, spoken in Nghe An and Thanh Hoa provinces in north-central Vietnam. Wikipedia gives 20,000 speakers (2002) and Joshua Project gives 25,000, so both figures are rough. Alternate names found are Tai Man Thanh, Tay Thanh, Thanh and Tai Then. The Wikipedia article notes some uncertainty about the name, which should be checked against other sources.

\section{Primary Educational Resources}
No primer, textbook or school programme for Tai Thanh turned up in search results.

\section{Linguistic and Reference Materials}
Wikipedia has a short article: \url{https://en.wikipedia.org/wiki/Tai_Thanh_language}. PHOIBLE appeared in the results with a source page, but I did not check whether it holds Tai Thanh data: \url{https://phoible.org/sources/82626}. The search summary also mentions Michel Ferlus (2008) on the Tai dialects of Nghe An (Tay Daeng, Tay Yo, Tay Muong), which is useful context for the wider group. I did not see a direct link to it and do not cite one. No Tai Thanh dictionary, grammar or orthography was found.

\section{Religious and Community Resources}
Joshua Project has a people-group profile: \url{https://joshuaproject.net/people_groups/15192}. It is written for missionary prayer, so its religious and livelihood details are secondary. It describes the people's own tribal, animistic religion. No scripture in Tai Thanh was found.

\section{Digital and Audio Resources}
Joshua Project has a language page for tmm: \url{https://joshuaproject.net/languages/tmm}. No audio or video was found.

\section{Curriculum Development Context}
The language is not known to be taught in schools, and almost nothing exists in the language. A curriculum would have to start from documentation: orthography, recorded texts and word lists. It would be developed with the community.
\end{document}